

% Template by Arnau, based on:
%%%%%%%%%%%%%%%%%%%%%%%%%%%%%%%%%%%%%%%%%
%
% Beamer Presentation
% LaTeX Template
% Version 1.0 (10/11/12)
% This template has been downloaded from:
% http://www.LaTeXTemplates.com
% License:
% CC BY-NC-SA 3.0 (http://creativecommons.org/licenses/by-nc-sa/3.0/)
%%%%%%%%%%%%%%%%%%%%%%%%%%%%%%%%%%%%%%%%%


%----------------------------------------------------------------------------------------
%	PACKAGES AND THEMES
%----------------------------------------------------------------------------------------

\documentclass[.08pt]{beamer}
% \usefonttheme{serif}
%% serif
%% professionalfonts
%% structurebold
%% structureitalicserif
%% structuresmallcapsserif

\mode<presentation> {

\usetheme{Boadilla}

% As well as themes, the Beamer class has a number of color themes
% for any slide theme. Uncomment each of these in turn to see how it
% changes the colors of your current slide theme.

% \usecolortheme{albatross}
%\usecolortheme{beaver}
%\usecolortheme{beetle}
%\usecolortheme{crane}
% \usecolortheme{dolphin}
%\usecolortheme{dove}
%\usecolortheme{fly}
%\usecolortheme{lily}
%\usecolortheme{orchid}
%\usecolortheme{rose}
%\usecolortheme{seagull}
\usecolortheme{seahorse}  % The them I usually use, because it resembles the PSI template.
% \usecolortheme{whale}
%\usecolortheme{wolverine}

%\setbeamertemplate{footline} % To remove the footer line in all slides uncomment this line
% \setbeamertemplate{footline}[page number] % To replace the footer line in all slides with a simple slide count uncomment this line
}

\setbeamertemplate{sections/subsections in toc}[sections numbered]  % Controls style of number in the table of contents.

% Packages:
\usepackage{adjustbox}
%\usepackage{algorithm}
%\usepackage{algpseudocode}
\usepackage{amsmath,amsfonts,amsthm,amssymb}
\usepackage{animate}
\usepackage{appendixnumberbeamer}
\usepackage{booktabs} % Allows the use of \toprule, \midrule and \bottomrule in tables
\usepackage{cancel}
%\usepackage{dutchcal}
\usepackage{float}
\usepackage{graphicx} % Allows including images
\usepackage{hhline}
\usepackage{listings}
\lstset{
  language=Python,
  basicstyle=\ttfamily,
  mathescape
}
%\usepackage{minted}
\usepackage{mathtools}
\usepackage{multicol}
\usepackage{multimedia}
\usepackage[authoryear]{natbib}
\usepackage{scrextend}
\changefontsizes{8pt}
%\usepackage{siunitx}
%\usepackage{soul}
\usepackage{tikz}
\usetikzlibrary{mindmap, trees, arrows, shapes, backgrounds, matrix, decorations.pathreplacing, decorations.pathmorphing, positioning, arrows.meta,plotmarks, calc}
\usepackage{verbatim}
\usepackage{xmpmulti}

\usepackage{hyperref}
\hypersetup{
	colorlinks=true,
	urlcolor=blue,
	linkcolor=blue
}

 %\bibliographystyle{unsrt}  % This style cites with numbers, e.g. [2]
%\bibliographystyle{unsrtnat}  % This style cites with names, e.g. [Adelmann, 2019]

\usepackage{caption}
\captionsetup{justification   = raggedright,
              singlelinecheck = false}


%%%%%%%%%%%%%%%%%%%%%%%%%%%%%%%%%
% Tikz stuff:
\usetikzlibrary{positioning}
\usetikzlibrary{patterns}
\usetikzlibrary{shapes}
\usetikzlibrary{matrix, arrows,decorations.pathmorphing}
\usetikzlibrary{decorations.markings}
\usetikzlibrary{fadings}
\usetikzlibrary{arrows.meta,bending}
%\usepackage{tikzscale}
%%%%%%%%%%%%%%%%%%%%%%%%%%%%%%%%%
%%%%%%%%%%%%%%%%%%%%%%%%%%%%%%%%%

\makeatletter
\setbeamertemplate{footline}{
  \leavevmode%
  \hbox{%
  \begin{beamercolorbox}[wd=.333333\paperwidth,ht=2.25ex,dp=1ex,center]{author in head/foot}%
    \usebeamerfont{author in head/foot}\insertshortauthor\expandafter\beamer@ifempty\expandafter{\beamer@shortinstitute}{}{~~(\insertshortinstitute)}
  \end{beamercolorbox}%
  \begin{beamercolorbox}[wd=.333333\paperwidth,ht=2.25ex,dp=1ex,center]{title in head/foot}%
    \usebeamerfont{title in head/foot}\let\hyperlink\@secondoftwo\insertshorttitle
  \end{beamercolorbox}%
  \begin{beamercolorbox}[wd=.333333\paperwidth,ht=2.25ex,dp=1ex,right]{date in head/foot}%
    \usebeamerfont{date in head/foot}\insertshortdate{}\hspace*{2em}
    \insertframenumber{} / \inserttotalframenumber\hspace*{2ex}
  \end{beamercolorbox}}%
  \vskip0pt%
}
\makeatother

\setbeamertemplate{navigation symbols}{}  % Remove navigation symbols
\setbeamersize{text margin left=1cm,text margin right=1cm}

\setbeamertemplate{caption}[numbered]

%% To-do lists
\usepackage{pifont}
\newcommand{\cmark}{\ding{51}}%
\newcommand{\xmark}{\ding{55}}%
\newcommand{\done}{\rlap{$\square$}{\raisebox{2pt}{\large\hspace{1pt}\cmark}}%
\hspace{-2.5pt}}
\newcommand{\wontfix}{\rlap{$\square$}{\large\hspace{1pt}\xmark}}

%% Maths definitions.
\DeclareMathOperator*{\argmin}{arg\,min}
\DeclareMathOperator*{\supp}{supp}
\DeclareMathOperator*{\Var}{\mathrm{Var}}
\DeclareMathOperator*{\Cov}{\mathrm{Cov}}
\DeclareMathOperator*{\E}{\mathbb{E}}
\DeclareMathOperator*{\MSE}{\text{MSE}}
\DeclareMathOperator*{\sgn}{\text{sgn}}
\DeclareMathOperator*{\Epsilon}{\mathcal{E}}
\DeclareMathOperator*{\bigO}{\mathcal{O}}
\DeclareMathOperator*{\smalls}{s}
\newcommand{\mvec}[1]{\boldsymbol{#1}}
\newcommand\underrel[2]{\mathrel{\mathop{#2}\limits_{#1}}}


\newcommand{\frametitlepsi}[1]{\frametitle{\hspace{0.8cm}\includegraphics[width=1.5cm]{logos/PSI.pdf}\hspace{1.1cm} #1}}

\title[\today]{IPPL Meeting}

\author{Tobia Claglüna}
\institute[LSM, PSI]{
    AMAS Group, LSM\\ % Your institution for the title page.
}
\date{\today}
\def \myEmail {tobia.clagluena@psi.ch}

\begin{document}
\def\vfilll{\vskip 0pt plus 1filll minus 0pt }
\begin{frame}
  % Title page image:
  \vspace{0.3cm}
  \begin{adjustbox}{width=\paperwidth, center}
    \begin{tikzpicture}
      \centering
      \filldraw[fill=lightgray!40!white, draw = none] (0,0) rectangle (3.3,0.5\textheight);
      % Logos:
      \node[anchor=south west,inner sep=0] at (0.4,3.4)
           {\includegraphics[width=2.5cm]{logos/PSI.pdf}};
      \node[anchor=south west,inner sep=0] at (0.5,1.8)
           {\includegraphics[width=2.4cm]{logos/eth_logo_kurz_pos-eps-converted-to.pdf}};
    \end{tikzpicture}
    \begin{tikzpicture}
      \node (heli) [anchor=south west,inner sep=0] at (0,0)
            {\includegraphics[height=0.5\textheight]{logos/PSI_helicopter}};
      \node [anchor=south west, inner sep = 1,fill=white, opacity=0.8] at (3.5, 0.0)
            {\tiny WIR SCHAFFEN WISSEN – HEUTE FÜR MORGEN};
    \end{tikzpicture}
    \begin{tikzpicture}
      \filldraw[fill=lightgray!40!white, draw = none] (0,0) rectangle (1,0.5\textheight);
    \end{tikzpicture}
  \end{adjustbox}
  \vspace{0.1cm}\\
  {\usebeamerfont{subtitle} \footnotesize \insertauthor\, ::  \insertinstitute}
  \vfilll
  \vspace{0.4cm}
  {\usebeamerfont{title} \LARGE \inserttitle}\\
  \vfilll
  \vspace{0.2cm}
  {\usebeamerfont{subtitle} \footnotesize \insertdate}\\
  \vfilll
  \null\hfill\tiny Contact: \url{\myEmail}
\end{frame}



\begin{frame}[c]{Friction Coefficient: DIH}
    \begin{figure}[!htb]
        \centering
        \captionsetup{justification=centering}
      \includegraphics[width=0.6\linewidth]{figures/friction_Nemittance.pdf}
      \caption{Poisson Solver possibly computes $ - F_d(\vec v)$?}
      \label{fig:Fd_DIH}
    \end{figure}
\end{frame}

\begin{frame}[c]{Diffusion Coefficient: DIH}
    \begin{enumerate}[$\cdot$]
        \item Simulation crashes after 200-400 timesteps.
        \item Computed $\boldsymbol D(\vec v)$ are positive definite and symmetric.
        \item Crashes in \texttt{scatterCIC()} as $\boldsymbol Q$ is too big (particles leave configuration space).
    \end{enumerate}
\end{frame}

\begin{frame}[c]{$\boldsymbol D$-Field Analysis: Maxwellian Test-Case}
    Diagonal elements of 3x3 matrices in the center of the domain $\Omega$, averaged over all velocity cells:
    \begin{multicols}{2}
    \begin{figure}[!htb]
        \centering
        \captionsetup{justification=centering}
      \includegraphics[width=1.03\linewidth]{figures/D_diagonal.pdf}
      \caption{Diagonal elements of $\Omega$ (without boundary cells).}
      \label{fig:Fd_asymptotic_behavior_2nd}
    \end{figure}
    \columnbreak
    \begin{figure}[!htb]
        \centering
        \captionsetup{justification=centering}
      \includegraphics[width=1.03\linewidth]{figures/reference_distribution_maxwell.png}
      \caption{Reference distribution \cite{Qiang_selfConsistentLangevin}}
      \label{fig:Fd_asymptotic_behavior_reference}
    \end{figure}
    \end{multicols}
\end{frame}

\begin{frame}[c]{$\boldsymbol D$-Field Analysis: Maxwellian Test-Case}
    Diagonal elements of 3x3 matrices averaged over all velocity cells:
    \begin{figure}[!htb]
        \centering
        \captionsetup{justification=centering}
      \includegraphics[width=1\linewidth]{figures/Dxx_avg_maxwellian.png}
      \caption{Diagonal elements of boundary cells $\partial \Omega$ compared to center cells $\Omega$.}
      \label{fig:Fd_asymptotic_behavior_2nd}
    \end{figure}
    $\implies$ Onesided Hessian is needed to resolve boundary cells properly.
\end{frame}

\begin{frame}[c]{$\boldsymbol D$-Field Analysis: Maxwellian Test-Case}
    Off-Diagonal elements $\boldsymbol D_{xy}$ of 3x3 matrices:
    \begin{figure}[!htb]
        \centering
        \captionsetup{justification=centering}
      \includegraphics[width=0.7\linewidth]{figures/Dxy_maxwellian_nolog.png}
      \caption{Distribution of off-diagonal entries on both the center domain and the boundary cells.}
      \label{fig:Fd_asymptotic_behavior_reference}
    \end{figure}
\end{frame}

\begin{frame}{Progress}
    \begin{itemize}
        \item[\done] Maxwellian Test-Case
        \item[$\square$] Onesided Hessian: $\boldsymbol D(\vec v) = \Gamma \frac{\partial^2}{\partial \vec v \partial \vec v}$
        \item[$\square$] Test how FD gradient impacts $F_d(\vec v)$ (instead of computing it in Fourier space)
        \item[$\square$] Check if \texttt{FFTPoissonSolver} solves for $\Delta \phi=-\rho$ or $\Delta \phi=\rho$?
        \item Analyse Diffusion Coefficients on Boundary Cells in DIH:
        \begin{itemize}
            \item[\done] Check that calculated $\boldsymbol D(\vec v)$ are positive definite and symmetric 
            \item[$\square$] Analyse Error introduced by Cholesky decomposition
        \end{itemize}
        \item Application to DIH:
        \begin{itemize}
            \item[\done] $F_d(\vec v)$
            \item[$\square$] $\boldsymbol D(\vec v)$
        \end{itemize}
    \end{itemize}
\end{frame}

\begin{frame}
\frametitle{(Adjusted) Timeline}

\begin{table}[]
    \def\arraystretch{1.5}
    \begin{tabular}{p{0.08\linewidth} | p{0.85\linewidth}}
        Date & Target Goals \\
        \hline \hline
        16/05 & Setup v-space datastructures in \texttt{LangevinParticles.hpp}. Add Friction coefficient. Add Solver for 2nd Rosenbluth potential $g(\vec v)$.\\
        23/05 & Analyse structure of $\boldsymbol D$. Finish Diffusion coefficient computation (via onesided Hessian operator). \\
        30/05 & Analyse interplay between collision coeff.`s (see whether Severin`s conclusions are confirmed or can be disproved). \\
              & Profiling of runtime and memory consumption. \\
        06/06 & Start improving most pressing bottlenecks. Start writing. \\
        17/07 & Submission. \\
    \end{tabular}
    \caption{Timeline with approximate milestones}
\end{table}

\end{frame}

\begin{frame}[allowframebreaks]
	\frametitle{References}
	\bibliographystyle{unsrtnat}
	\bibliography{references.bib}
\end{frame}

\end{document}
